% Neuausrichtung der Masterarbeit ab 12.09.2026
% Vorheriger Stand: Archive/Thesis_PreRealignment_2026-09-12/Kapitel/

\section{Diskussion und Beantwortung der Forschungsfragen}
\label{sec:discussion}

\subsection{Einordnung der zentralen Ergebnisse}
\label{sec:discussion_overview}

Die Ergebnisse der prototypischen Untersuchung zeigen, dass die zentrale
Herausforderung einer KI-gestützten Überführung heterogener
Engineering-Information in SysML~v2 nicht primär in der Erzeugung syntaktisch
plausibler Modellnotation liegt. Entscheidend ist vielmehr, wie aus
heterogenen und teilweise unvollständigen Ausgangsinformationen ein
kontrollierter, nachvollziehbarer und autorisierter Engineering-Zustand
entsteht, der anschließend in ein Zielmodell überführt werden kann. Diese
Beobachtung entspricht der in der Literatur beschriebenen Herausforderung,
LLM-basierte Verfahren trotz probabilistischer Ausgabe, möglicher
Fehlinterpretationen und begrenzter Verlässlichkeit in kontrollierte
Engineering-Prozesse einzubetten \cite{Topcu2025,Barke2023,Sergeyuk2025}.

Die Evaluation macht dabei deutlich, dass ein direkter Verarbeitungspfad von
einem Legacy-Dokument zu einem generierten SysML-v2-Artefakt für den
untersuchten Anwendungsfall nicht ausreicht. Während der formativen
End-to-End-Erprobung mussten mehrere zunächst zusammengefasste
Verantwortlichkeiten explizit getrennt werden. Dies betrifft insbesondere die
Trennung von Engineering Evidence und zusätzlichem Verarbeitungskontext, die
Bildung gemeinsamer Engineering Subjects vor deren perspektivischer
Interpretation sowie die Unterscheidung von fachlicher Bedeutung,
Modellplatzierung und zielmodellspezifischer Formulierung. Der resultierende
Informationsfluss ist daher weniger als lineare Transformationskette, sondern
als Folge kontrollierter Zustands- und Authority-Übergänge zu verstehen.

Eine zweite zentrale Beobachtung betrifft die Rolle probabilistischer
Interpretation. Die im Referenzlauf beobachteten unterschiedlichen
Persona-Klassifikationen zeigen, dass mehrere KI-basierte Perspektiven nicht
notwendigerweise zu einer eindeutigen fachlichen Aussage konvergieren. Im
entwickelten Architekturkonzept wird diese Varianz deshalb nicht als
automatischer Abstimmungsmechanismus verwendet. Stattdessen wird sie als
Review-Evidence persistiert und einem expliziten Human Engineering Review
zugeführt. Dies steht im Einklang mit Ansätzen des Human-in-the-Loop, bei
denen menschliche Entscheidungskompetenz gezielt an solchen Stellen erhalten
bleibt, an denen automatisierte Verarbeitung keine hinreichende
Entscheidungsautorität besitzt \cite{MosqueiraRey2023,Lazaros2026}. Die
Ergebnisse des Prototyps zeigen zugleich, dass diese Human Authority nicht nur
an einer einzigen Stelle benötigt wird, sondern in unterschiedlichen
Verantwortungsbereichen des Informations- und Modellbildungsprozesses.

Als dritter Befund ergibt sich, dass fachlich autorisierte Engineering
Information und deren konkrete Modellrepräsentation nicht identisch sind.
Eine freigegebene Engineering-Aussage beantwortet zunächst, welcher fachliche
Inhalt für die weitere Verarbeitung autorisiert ist. Davon getrennt muss
entschieden werden, an welcher Stelle und in welcher Form dieser Inhalt im
Zielmodell repräsentiert wird. Erst anschließend kann die konkrete
SysML-v2-Formulierung erzeugt und technisch geprüft werden. Die im
Referenzlauf getrennt persistierten Engineering-, Placement-, Formulation- und
Validation-Zustände zeigen diese unterschiedlichen Verantwortungsbereiche
auch in der prototypischen Realisierung. Die formale Ausdrucksfähigkeit von
SysML~v2 und textbasierte Engineering- beziehungsweise
Architecture-as-Code-Ansätze stellen hierfür die technische Grundlage bereit,
ersetzen jedoch nicht die vorgelagerten fachlichen Entscheidungen
\cite{SysMLv2,Stirbu2022,Bucaioni2025}.

Ein weiterer zentraler Befund betrifft die Nachvollziehbarkeit des
Entstehungsprozesses. Das Ergebnis des Referenzlaufs besteht nicht lediglich
aus einem generierten SysML-v2-Modell. Parallel zum Modell wurde ein
persistierter Evidence-, Decision-, Provenance- und
Traceability-Artefaktsatz erzeugt, über den Source-Aussagen,
maschinelle Zwischenzustände, Human Decisions, Modellzustände,
Validierungsergebnisse und der veröffentlichte Output miteinander verbunden
bleiben. Damit wird Provenance nicht erst nach Abschluss der Modellbildung
ergänzt, sondern während der Verarbeitung fortlaufend aufgebaut. Die
grundsätzliche Bedeutung expliziter Provenance für die Rekonstruktion von
Daten- und Verarbeitungshistorien ist sowohl in allgemeinen
Provenance-Modellen als auch in Engineering- und
Traceability-Kontexten etabliert \cite{MoreauMissier2013,Simmhan2005,ISO29148}.

Die Multi-Source-Erprobung erweitert diese Betrachtung um die
projektübergreifende beziehungsweise projektweite Synthese mehrerer
Engineering Sources. Dabei zeigte sich, dass source-lokale Provenance und
projektweite Modellbildung nicht als gegensätzliche Anforderungen behandelt
werden müssen. Die einzelnen Quellen können zunächst unabhängig verarbeitet
und autorisiert werden, während ihre freigegebenen Informationszustände
anschließend gemeinsam in die Modellbildung eingehen. Voraussetzung ist,
dass Source-, Processing- und Authority-Bezüge auch nach dem Übergang in die
projektweite Verarbeitung erhalten bleiben. Die Evaluation bezieht sich dabei
auf die im Testfall registrierten und zugelassenen Projektinformationen und
begründet keine Aussage über die Vollständigkeit beliebiger realer
Engineering-Datenbestände.

\paragraph{Umgang mit Findings und Architekturkorrekturen}

Die formative Evaluation hatte innerhalb der Entwicklung damit nicht nur die
Funktion einer nachgelagerten Prüfung. Beobachtete Abweichungen wurden auch als
Rückkopplung für die Weiterentwicklung des Artefakts verwendet. Dieses Vorgehen
entspricht dem iterativen Charakter gestaltungsorientierter Forschung, bei der
Evaluationsergebnisse zur begründeten Überarbeitung des entwickelten Artefakts
herangezogen werden können \cite{Wieringa2014}. Dabei führte jedoch nicht jedes
Finding automatisch zu einer Änderung des Architekturkonzepts.

Entscheidend war vielmehr die Ursache des Befunds. Zeigte ein Finding, dass eine
angenommene Verantwortungs-, Informations-, Provenance- oder Authority-Grenze
konzeptionell nicht ausreichte, wurde die daraus abgeleitete Korrektur als
Architekturentscheidung formalisiert und in einem Architecture Decision Record
dokumentiert. Die korrigierte Architektur wurde anschließend implementiert und
am betroffenen Verarbeitungspfad erneut geprüft. Die ursprüngliche Beobachtung
blieb dabei als historische Evaluationsevidenz erhalten und wurde durch den
späteren erfolgreichen Retest nicht rückwirkend ersetzt. Beispiele hierfür sind
die Einführung source-grounded Evidence vor der Persona-Interpretation sowie
die Präzisierung der source-lokalen und projektweiten Authority-Grenzen der
Multi-Source-Verarbeitung.

Davon getrennt wurden technische Implementierungsfehler behandelt, die den
Verifikationspfad blockierten, ohne eine Änderung der konzeptionellen
Architektur zu erfordern. Solche Fehler wurden korrigiert und retestet, jedoch
nicht als eigenständige Architekturentscheidung interpretiert.

Nach Erreichen des für die Arbeit vorgesehenen und erfolgreich durchlaufenen
Proof-of-Concept-Umfangs wurde die prototypische Realisierung nicht weiter in
Richtung eines produktionsreifen Systems optimiert. Verbleibende nicht
blockierende Findings, insbesondere zur Bedienbarkeit, Laufzeit,
Relationship-Konsolidierung, Lifecycle-Unterstützung und zur noch begrenzten
Abdeckung von SysML-v2-Konstrukten, wurden deshalb bewusst nicht vollständig
aufgelöst. Sie bleiben im vollständigen Findings-Register erhalten und werden
als Limitationen beziehungsweise als Ansatzpunkte für die weitere Entwicklung
in den Ausblick übernommen.

Diese Trennung ist für die Interpretation der Ergebnisse relevant. Ein
erfolgreicher finaler End-to-End-Lauf bedeutet nicht, dass sämtliche während
der formativen Evaluation beobachteten Einschränkungen beseitigt wurden.
Er zeigt vielmehr, dass der definierte und architektonisch konsolidierte
Proof-of-Concept-Pfad unter den in dieser Arbeit untersuchten Bedingungen
vollständig ausgeführt werden konnte.

Aus den Ergebnissen lassen sich damit vier miteinander verbundene
Diskussionsschwerpunkte ableiten: die Kontrolle der Informationsbasis, die
Rolle von Interpretationsvarianz und Human Authority, die Trennung von
Engineering Meaning und Zielmodellbildung sowie die durchgängige Persistenz
von Evidence, Provenance und Traceability. Ergänzt wird diese Betrachtung um
die Frage, wie mehrere source-lokale Informationspfade kontrolliert in einen
gemeinsamen projektweiten Modellzustand überführt werden können. Diese
Aspekte werden in den folgenden Abschnitten einzeln diskutiert, bevor daraus
in Abschnitt~\ref{sec:discussion_research_questions} die Forschungsfragen
beantwortet werden.

\subsection{Kontrollierte Informationsbasis und source-grounded Evidence}
\label{sec:discussion_information_basis}

Die Ergebnisse der formativen Erprobung zeigen, dass die Zuverlässigkeit der
KI-gestützten Verarbeitung nicht allein von der technischen Lesbarkeit oder
formalen Struktur der Eingangsdaten abhängt. Bereits vor der eigentlichen
Interpretation muss festgelegt sein, welche Informationen innerhalb eines
Projekts überhaupt fachliche Aussagekraft besitzen und in welcher Rolle sie
verwendet werden dürfen. Besonders deutlich wurde dies durch den Befund, dass
Orchestrierungs-, Task- und Prozessinformationen in einem frühen
Verarbeitungsstand teilweise als positive Engineering Information
interpretiert wurden. Der Inhalt war für das LLM verfügbar, besaß jedoch nicht
dieselbe fachliche Authority wie die eigentliche Engineering Source.

Für das entwickelte Architekturkonzept folgt daraus eine Trennung zwischen
\emph{Engineering Sources} und \emph{Context Sources}. Engineering Sources
bilden die mögliche fachliche Evidenzbasis für Aussagen über das betrachtete
System. Context Sources können dagegen Terminologie, Klassifikation,
Interpretation oder projektbezogene Entscheidungen unterstützen, dürfen jedoch
nicht eigenständig zur Begründung einer Engineering-Aussage verwendet werden.
Die entscheidende Randbedingung besteht damit nicht darin, zusätzliches Wissen
grundsätzlich aus der Verarbeitung auszuschließen, sondern die Rollen
unterschiedlicher Informationsquellen explizit zu kontrollieren.

Diese Trennung ist insbesondere für probabilistische Sprachmodelle relevant.
LLM-basierte Systeme können sprachlich plausible Engineering-Artefakte
erzeugen, ohne dass aus der Plausibilität unmittelbar auf die fachliche
Korrektheit oder die Herkunft der enthaltenen Aussagen geschlossen werden
kann. Untersuchungen zur Verwendung generativer Modelle in Engineering- und
Softwarekontexten weisen entsprechend auf Fehlerbilder und auf die begrenzte
Verlässlichkeit scheinbar überzeugender Modelloutputs hin
\cite{Topcu2025,Barke2023,Sergeyuk2025}. Im hier untersuchten Ansatz wird
dieses Problem nicht durch den Anspruch gelöst, die probabilistische
Interpretation selbst fehlerfrei zu machen. Stattdessen wird ihre mögliche
fachliche Wirkung an eine kontrollierte Evidenzbasis gebunden.

Die hierfür eingeführte source-grounded Evidence bildet eine explizite
Zwischenschicht zwischen Source und Interpretation. Ein Evidence Item hält
nicht nur eine extrahierte Aussage fest, sondern bindet diese an die
registrierte Source und die zugehörige Position beziehungsweise
Source Projection. Erst auf dieser Grundlage werden Canonical Engineering
Subjects gebildet und anschließend durch die unterschiedlichen
Persona-Perspektiven interpretiert. Die formative Evaluation zeigte, dass die
frühere Persona-first-Verarbeitung zu voneinander abweichenden
Subject-Populationen führte. Durch die Vorverlagerung einer gemeinsamen
Evidence- und Subject-Identität beziehen sich unterschiedliche
Interpretationen dagegen auf denselben fachlichen Referenzgegenstand.

Damit erhält Provenance bereits am Beginn des Engineering Information Flow
eine operative Funktion. Provenance beschreibt grundsätzlich die Herkunft
eines Datenobjekts und die Verarbeitungsschritte, durch die ein bestimmter
Zustand entstanden ist \cite{MoreauMissier2013,Simmhan2005}. Im entwickelten
Ansatz wird diese Herkunft nicht erst nachträglich zum erzeugten Modell
dokumentiert. Source-, Evidence- und Processing-Bezüge werden bereits beim
Entstehen der jeweiligen Zwischenzustände persistiert und über die
nachfolgenden Authority-Grenzen weitergeführt. Dies schafft gleichzeitig die
Grundlage für die in Abschnitt~\ref{sec:discussion_traceability} diskutierte
End-to-End-Traceability.

Die Ergebnisse zeigen außerdem, dass die Anforderungen an die Eingangsdaten
nicht mit einer Forderung nach vollständigen oder bereits widerspruchsfreien
Legacy-Dokumenten gleichgesetzt werden sollten. Gerade Brownfield- und
Legacy-Kontexte sind dadurch gekennzeichnet, dass relevante
Engineering-Information über bestehende Dokumente und unterschiedliche
Darstellungsformen verteilt sein kann. MBSE verfolgt demgegenüber das Ziel,
Engineering-Information konsistenter in miteinander verbundenen
Modellstrukturen zu erfassen und nutzbar zu machen \cite{Hendriks2022}.
Für die Migration ergibt sich daraus jedoch nicht, dass fehlende Information
durch die KI ergänzt werden darf.

Der entwickelte Verarbeitungspfad unterscheidet deshalb zwischen
\emph{nicht eindeutig interpretierbarer} und \emph{nicht vorhandener}
Information. Unsicherheit über eine vorhandene Source-Aussage kann als
Interpretationsvarianz oder Open Question in den Human Review gelangen.
Information, für die keine ausreichende Evidence in den zugelassenen Sources
vorliegt, erhält dagegen keine positive Engineering Authority. Der
fail-closed Ansatz verhindert dabei nicht jede fehlerhafte Interpretation,
begrenzt jedoch deren automatische Weitergabe in einen autorisierten
Engineering- oder Modellzustand.

Eine weitere Konsequenz betrifft die Bedeutung von Vollständigkeit.
Die Evaluation kann nur zeigen, dass die für einen Verarbeitungslauf
registrierten und zugelassenen Sources entsprechend der implementierten
Verträge verarbeitet und ihre resultierenden Informationszustände
nachvollziehbar weitergeführt werden. Daraus folgt keine Aussage, dass die
registrierten Sources sämtliche für das reale System relevanten Informationen
enthalten. Diese Grenze ist insbesondere für eine spätere industrielle
Anwendung wesentlich: Die Architektur kann die kontrollierte Verarbeitung
einer definierten Informationsbasis unterstützen, die Auswahl und fachliche
Vollständigkeit dieser Informationsbasis bleibt jedoch eine vorgelagerte
Engineering-Verantwortung.

Auch die Anforderungen des Requirements Engineering an nachvollziehbare
Herkunft und Beziehungen zwischen Engineering-Information sprechen dafür,
Traceability nicht erst am finalen Modell anzusetzen
\cite{ISO29148}. Für die KI-gestützte Migration bedeutet dies, dass
Eingangsdaten neben ihrer technischen Verarbeitbarkeit insbesondere eine
stabile Source-Identität, eine explizite Informationsrolle und eine
nachvollziehbare Provenance benötigen. Erst diese Randbedingungen ermöglichen
es, nachgelagerte KI-Interpretationen von der fachlichen Evidenzbasis zu
unterscheiden.

Damit verschiebt sich die Frage nach geeigneten Eingangsdaten von einer rein
formatbezogenen Betrachtung hin zu einer Governance-Frage. Für den
untersuchten Ansatz ist weniger entscheidend, ob ein Legacy-Dokument bereits
einer bestimmten Modellierungsnotation folgt. Entscheidend ist, ob es
registriert, eindeutig identifizierbar, hinsichtlich seiner Informationsrolle
klassifiziert und auf konkrete fachliche Aussagen zurückführbar ist. Auf
dieser kontrollierten Informationsbasis können anschließend probabilistische
Interpretation, Human Review und Modellbildung erfolgen. Die Konsequenzen für
die formale Beantwortung der ersten Teilforschungsfrage werden in
Abschnitt~\ref{sec:discussion_research_questions} zusammengeführt.

\subsection{Interpretationsvarianz und Human Authority}
\label{sec:discussion_human_authority}

Die Ergebnisse der formativen Evaluation zeigen, dass mehrere
KI-Perspektiven nicht automatisch zu einer eindeutigen oder fachlich richtigen
Engineering-Aussage führen. Unterschiedliche Personas können denselben
source-grounded Gegenstand verschieden klassifizieren, während in anderen
Fällen mehrere Perspektiven dieselbe Unsicherheit teilen. Für das entwickelte
Architekturkonzept folgt daraus, dass Konsens und Varianz zunächst als
Information über den Interpretationszustand behandelt werden und nicht als
Engineering Authority.

Multi-Agent-Ansätze werden in der Literatur unter anderem mit dem Ziel
untersucht, durch den Austausch mehrerer Modellinstanzen die Qualität von
Reasoning und faktischen Aussagen zu verbessern. Du et al. zeigen
beispielsweise, dass Multi-Agent Debate unter bestimmten Bedingungen zu
verbesserten Ergebnissen führen kann \cite{Du2024}. Eine solche
Mehrheits- oder Debattenlogik ist jedoch nicht mit einer fachlichen
Freigabeentscheidung im Systems Engineering gleichzusetzen. Auch eine
übereinstimmende Aussage mehrerer probabilistischer Instanzen besitzt zunächst
nur den Status einer maschinell erzeugten Interpretation.

Der Prototyp verwendet mehrere Persona-Perspektiven deshalb nicht als
Abstimmungsverfahren zur automatischen Ermittlung einer vermeintlich richtigen
Antwort. Ihre Funktion besteht vielmehr darin, unterschiedliche plausible
Interpretationen desselben zuvor identifizierten Engineering Subjects
sichtbar zu machen. Voraussetzung dafür ist die in
Abschnitt~\ref{sec:discussion_information_basis} diskutierte gemeinsame
Evidence- und Subject-Identität. Erst wenn die Personas denselben
Referenzgegenstand interpretieren, kann eine Abweichung ihrer Ergebnisse
sinnvoll als Interpretationsvarianz betrachtet werden.

Diese Trennung entstand unmittelbar aus der formativen Evaluation. In einem
früheren Architekturstand erzeugten die Persona-Perspektiven teilweise eigene
Subject-Populationen. Unterschiede zwischen den Ergebnissen vermischten damit
zwei Fragestellungen: ob überhaupt derselbe Engineering-Gegenstand erkannt
worden war und wie dieser Gegenstand interpretiert werden sollte. Die
Architekturkorrektur verlagerte deshalb source-grounded Evidence Detection und
Canonical Subject Discovery vor die Persona-Interpretation. Der entsprechende
Befund wurde als Architekturentscheidung in die weitere Systemgestaltung
überführt und nach der Umsetzung erneut geprüft
(vgl. Anhang~\ref{app:findings_register}).

Der Referenzlauf zeigt zusätzlich, dass zwischen
\emph{Interpretationsvarianz} und \emph{gemeinsamer Ambiguität} unterschieden
werden muss. Weichen die Persona-Ergebnisse voneinander ab, liegt eine
beobachtbare Varianz zwischen den Perspektiven vor. Liefern dagegen mehrere
Perspektiven dieselben alternativen Möglichkeiten oder weisen sie auf dieselbe
fehlende Information hin, verschwindet die Unsicherheit nicht dadurch, dass
kein Inter-Persona-Konflikt besteht. Consensus darf daher nicht grundsätzlich
mit hoher Sicherheit gleichgesetzt werden. Diese Unterscheidung war ebenfalls
Gegenstand der formativen Findings und bleibt für die Gestaltung zukünftiger
Review-Unterstützung relevant.

Damit verändert sich auch die Funktion des Human-in-the-Loop. Der Mensch wird
nicht lediglich als nachgelagerter Korrekturmechanismus eingesetzt, wenn ein
LLM sichtbar scheitert. Vielmehr besitzt er an definierten Stellen eine
explizite Authority über den Übergang in einen autorisierten
Engineering-Zustand. Human-in-the-Loop-Ansätze betonen allgemein die
Einbindung menschlicher Expertise in maschinelle Lern- und
Entscheidungsprozesse, insbesondere wenn automatische Ergebnisse überprüft,
korrigiert oder ergänzt werden müssen
\cite{MosqueiraRey2023,Lazaros2026}. Auch Untersuchungen zu menschlicher
Intervention in Multi-Agent-LLM-Verfahren zeigen, dass die Rolle des Menschen
nicht durch die Interaktion mehrerer Modelle vollständig entfällt
\cite{TriemDing2024}.

Die Evaluation des Prototyps führte dabei zu einer weiteren Präzisierung:
Human Authority ist nicht als einzelnes finales Approval Gate ausreichend.
Die fachliche Autorisierung einer Engineering-Aussage, die Entscheidung über
ihre Modellplatzierung, eine gegebenenfalls erforderliche
Target-Model-Formulation sowie die abschließende Freigabe des erzeugten
Modellartefakts betreffen unterschiedliche Verantwortungsbereiche. Deshalb
wurden diese Entscheidungen im finalen Architekturstand als getrennte
Authority-Grenzen realisiert.

Besonders deutlich wird dies bei semantischen Relationships. Die formative
Erprobung zeigte, dass die Autorisierung einzelner Subjects nicht automatisch
die zwischen ihnen vorgeschlagenen fachlichen Beziehungen einschließt.
Akzeptierte Relationships mussten deshalb ebenfalls explizit Bestandteil der
autorisierten Engineering Information werden. Umgekehrt darf eine
nachgelagerte Modellbildungsstufe eine noch ungelöste fachliche Beziehung
nicht durch eine erneute LLM-Interpretation implizit entscheiden. Wo eine
eindeutige Weiterverarbeitung aus den bereits autorisierten Zuständen nicht
möglich ist, bleibt der Zustand unresolved oder wird erneut einer zuständigen
Human Authority vorgelegt.

Damit erfüllt das fail-closed Verhalten eine andere Funktion als die
Vermeidung sämtlicher KI-Fehler. Es verhindert nicht, dass eine Persona eine
unpassende Klassifikation vorschlägt. Es begrenzt jedoch, unter welchen
Bedingungen eine solche Interpretation in einen autorisierten
Engineering- oder Modellzustand überführt werden kann. Der Human Review bildet
dadurch keine nachträgliche Qualitätskontrolle eines bereits feststehenden
Modells, sondern einen Bestandteil des eigentlichen
Informationsverarbeitungsprozesses.

Gleichzeitig zeigte die formative Untersuchung Grenzen dieses Ansatzes. Die
Review-Oberflächen erzeugten teilweise hohen Lese- und Interaktionsaufwand,
gemeinsame Ambiguität wurde noch nicht in allen Fällen optimal von echter
Persona-Varianz dargestellt und einzelne Lifecycle-Fälle erfordern weiterhin
explizite manuelle Eingriffe. Diese verbleibenden, nicht blockierenden
Findings wurden nach Abschluss des vorgesehenen Proof-of-Concept-Umfangs nicht
weiter optimiert und werden in den Limitationen sowie im Ausblick erneut
aufgegriffen (vgl. Anhang~\ref{app:findings_register}).

Für den untersuchten Ansatz ergibt sich damit eine klare Rollenverteilung:
Probabilistische KI unterstützt die Erzeugung und Gegenüberstellung
möglicher Interpretationen; Consensus und Variance strukturieren die daraus
entstehende Review-Evidence. Die fachliche Authority verbleibt jedoch an
explizit definierten Human-Decision-Grenzen. Die Konsequenzen dieses
Ergebnisses für die zweite Teilforschungsfrage werden in
Abschnitt~\ref{sec:discussion_research_questions} zusammengeführt.

\subsection{Von Engineering Meaning zur SysML-v2-Repräsentation}
\label{sec:discussion_model_formation}

Die formative Evaluation zeigte, dass die Überführung fachlich autorisierter
Engineering Information in SysML~v2 nicht als unmittelbare Übersetzung
verstanden werden kann. Eine fachlich akzeptierte Aussage legt zunächst fest,
\emph{was} über das betrachtete System als autorisiert gilt. Daraus folgt
jedoch nicht eindeutig, \emph{wo} diese Information im Zielmodell einzuordnen
ist, \emph{durch welchen Modellkonstrukt} sie repräsentiert werden soll oder
\emph{wie} die resultierende SysML-v2-Formulierung konkret auszusehen hat.
Gerade diese zunächst nicht ausreichend getrennten Verantwortlichkeiten
führten während der End-to-End-Erprobung zu mehreren architekturrelevanten
Findings (vgl. Anhang~\ref{app:findings_register}).

Für den finalen Architekturstand wurden deshalb vier aufeinanderfolgende
Verantwortungsbereiche unterschieden:

\begin{center}
Engineering Meaning\\
$\downarrow$\\
Model Placement / Representation\\
$\downarrow$\\
Target-Model Formulation\\
$\downarrow$\\
deterministische SysML-v2-Serialisierung.
\end{center}

\emph{Engineering Meaning} bezeichnet dabei den fachlich autorisierten
Informationszustand. Die nachfolgende Placement- beziehungsweise
Representation-Entscheidung bestimmt, an welcher Stelle und als welcher
zugelassene Modelltyp dieser Inhalt im Zielmodell repräsentiert werden soll.
Die Target-Model Formulation konkretisiert anschließend die
zielspezifische Ausprägung innerhalb der zulässigen Modellierungsnotation.
Erst danach erfolgt die deterministische Serialisierung in ein
SysML-v2-Artefakt.

Diese Trennung präzisiert zugleich die Rolle ontologischer und
modellmethodischer Leitplanken. Ontologien können Begriffe, Kategorien und
semantische Beziehungen formalisieren und damit die konsistente Interpretation
und Integration von Engineering-Information unterstützen
\cite{OrellanaMandrick2019,Lu2022,Kulvatunyou2022}. Auch standardisierte
Top-Level-Ontologien zeigen den grundsätzlichen Nutzen expliziter
Kategorisierungs- und Relationsregeln für interoperable semantische Modelle
\cite{ISO21838_2_2021}. Für den hier untersuchten Transformationsprozess
reichen solche Regeln allein jedoch nicht aus.

Die formative Erprobung zeigte insbesondere, dass eine fachliche
Klassifikation nicht zwangsläufig einen eindeutigen SysML-v2-Konstrukt
vorschreibt. Ein Engineering Information Type beschreibt zunächst die
fachliche Bedeutung einer Information. Die konkrete Zielrepräsentation hängt
zusätzlich vom verwendeten Modellframework, dem Modellkontext und den für den
Proof of Concept zugelassenen Modellierungsmustern ab. Umgekehrt darf die
Verfügbarkeit eines technisch passenden SysML-v2-Konstrukts nicht darüber
entscheiden, ob die zugrunde liegende Engineering-Aussage fachlich autorisiert
ist.

Damit erhält die Ontologie im entwickelten Konzept die Funktion einer
\emph{Leitplanke}, nicht die einer vollständigen Entscheidungsinstanz. Sie
begrenzt zulässige Interpretationen und unterstützt die semantische
Konsistenz. Die Authority über Engineering Meaning sowie über nicht eindeutig
ableitbare Placement- oder Formulation-Entscheidungen verbleibt dagegen an
den jeweils vorgesehenen Human-Decision-Grenzen.

Diese Abgrenzung wurde während der Evaluation besonders deutlich, als der
frühere Modellbildungsprozess für bereits akzeptierte Relationships erneut
LLM-basierte Placement-Perspektiven aufrief. Dadurch entstand innerhalb der
Model Assembly eine neue semantische Entscheidungsinstanz, obwohl die
Assembly eigentlich einen bereits autorisierten Informationszustand
materialisieren sollte. Die daraufhin vorgenommene Architekturkorrektur
begrenzte die Model Assembly auf deterministisch ausführbare Übergänge. Die
entsprechende Beobachtung, Architekturkorrektur und der erfolgreiche Retest
sind im Findings-Register dokumentiert
(vgl. Anhang~\ref{app:findings_register}).

Diese Grenze besitzt eine wichtige Konsequenz für unresolved States. Kann eine
fachlich autorisierte Information mit den verfügbaren Modellierungsregeln
nicht eindeutig abgebildet werden, darf die Assembly keine neue semantische
Interpretation erzeugen, nur um einen technisch vollständigen Modellzustand
zu erreichen. Der Zustand bleibt stattdessen explizit unresolved oder wird
einer geeigneten Human Authority vorgelegt. Fail-closed bedeutet in diesem
Zusammenhang somit, dass fehlende Modellierungsentscheidungen nicht durch
implizite probabilistische Entscheidungen verdeckt werden.

Auch die Einführung einer expliziten Target-Model-Formulation-Stufe entstand
aus der formativen Untersuchung. Der ursprüngliche Pfad führte zu direkt von
fachlich autorisierter Information und Placement in Richtung
SysML-v2-Generierung. Im weiteren Verlauf zeigte sich jedoch, dass auch nach
einer geklärten Platzierung noch eine eigenständige Entscheidung über die
konkrete modellsprachliche Ausprägung erforderlich sein kann. Dieser
Verantwortungsbereich wurde deshalb vor die deterministische Serialisierung
gezogen und als eigener persistierter Authority-Zustand realisiert. Der
Referenzlauf~930444 zeigt mit der Target-Model-Formulation des
\emph{Microscope Operator} exemplarisch eine solche explizite Entscheidung
(vgl. Abschnitt~\ref{sec:results_traceability}).

SysML~v2 unterstützt durch seine formale Sprachdefinition und textuelle
Notation eine solche Trennung zwischen internem Engineering-Zustand und
anschließender technischer Repräsentation
\cite{SysMLv2,Teodorov2025}. Textbasierte Engineering- und
Architecture-as-Code-Ansätze ermöglichen darüber hinaus, Modellartefakte
deterministisch zu erzeugen, versionierbar zu behandeln und in automatisierte
Werkzeugketten einzubetten \cite{Stirbu2022,Bucaioni2025}. Die technische
Erzeugbarkeit eines Artefakts beantwortet jedoch weiterhin nicht die Frage,
ob die darin enthaltene Engineering-Aussage fachlich richtig oder ausreichend
begründet ist.

Dies gilt ebenfalls für die anschließende technische Validierung. Der im
Referenzlauf erfolgreich durch SYSIDE verarbeitete Output zeigt, dass das
erzeugte Artefakt innerhalb der verwendeten Werkzeug- und Sprachumgebung
technisch verarbeitbar war. Die Validierung ersetzt weder die vorausgehende
Engineering Authority noch die Model Authority. Entsprechend bleiben
Generierung, technische Validation und abschließende Human Release im
Architekturkonzept getrennte Zustände.

Die Ergebnisse begrenzen zugleich die Reichweite der entwickelten
Modellierungsleitplanken. Der Proof of Concept implementiert einen definierten
Ausschnitt der relevanten SysML-v2-Konstrukte. Das Findings-Register weist
weiterhin eine unvollständige Abdeckung unter anderem spezifischer Ports,
Interfaces, Parts, Actions, States und Occurrences aus. Diese Einschränkung
wurde nach Abschluss des vorgesehenen Proof-of-Concept-Umfangs nicht durch
eine Erweiterung der Construct-Bibliothek aufgelöst, sondern als offene
Systemgrenze und Gegenstand des Ausblicks erhalten
(vgl. Anhang~\ref{app:findings_register}). Der erfolgreiche Referenzlauf ist
daher als Nachweis des implementierten Modellierungsausschnitts und nicht als
Nachweis einer vollständigen SysML-v2-Abdeckung zu verstehen.

Für die dritte Teilforschungsfrage ergibt sich daraus bereits eine wesentliche
Präzisierung: Ontologische Regeln sind ein notwendiger Bestandteil der
Leitplanken, reichen für eine kontrollierte Überführung in modulare
SysML-v2-Strukturen jedoch nicht allein aus. Sie müssen durch Regeln für
zulässige Zielrepräsentationen, explizite Placement- und
Formulation-Entscheidungen sowie deterministische Validierungs- und
Serialisierungsgrenzen ergänzt werden. Die formale Antwort auf die
Teilforschungsfrage wird in
Abschnitt~\ref{sec:discussion_research_questions} zusammengeführt.

\subsection{Persistierte Evidence, Provenance und End-to-End-Traceability}
\label{sec:discussion_traceability}

Die Ergebnisse des Referenzlaufs zeigen, dass der kontrollierte
Engineering-Prozess nicht allein durch das resultierende SysML-v2-Modell
beschrieben werden kann. Zusätzlich entsteht ein persistierter Evidence-,
Decision-, Provenance- und Traceability-Artefaktsatz, über den nachvollziehbar
bleibt, aus welchen Eingangsinformationen ein Modellzustand entstanden ist,
welche maschinellen Zwischenzustände vorlagen und an welchen Stellen
menschliche Entscheidungen den weiteren Verarbeitungspfad autorisiert haben.
Der exemplarische Informationspfad in
Abschnitt~\ref{sec:results_traceability} zeigt diese Verkettung vom
Source-Inhalt bis zum publizierten Modellelement.

Provenance beschreibt allgemein die Herkunft von Daten sowie die Prozesse und
Transformationen, durch die ein bestimmter Zustand entstanden ist
\cite{MoreauMissier2013,Simmhan2005}. Für den hier untersuchten
KI-gestützten Engineering-Prozess reicht es jedoch nicht aus, lediglich die
ursprüngliche Source eines finalen Modellelements zu kennen. Zwischen Source
und Modell liegen probabilistische Interpretationen, Human Decisions,
Modellierungsentscheidungen und technische Transformationen. Die Herkunft
eines finalen Artefakts umfasst deshalb neben der Source-Provenance auch die
jeweiligen Processing- und Authority-Zustände.

Diese Unterscheidung wird insbesondere bei KI-generierten Zwischenzuständen
relevant. Eine LLM-Antwort besitzt im entwickelten Ansatz zunächst keine
Engineering Authority. Sie wird als Verarbeitungsergebnis beziehungsweise
Review-Evidence persistiert und bleibt von einem später autorisierten Zustand
unterscheidbar. Im Referenzlauf ist dies beispielsweise an den
Persona-Interpretationen des Subjects \emph{Microscope Workstation}
nachvollziehbar. Die voneinander abweichenden Klassifikationen bleiben
erhalten, obwohl der nachfolgende Human Engineering Review einen konkreten
Information Type autorisiert. Der Human Decision State ersetzt damit nicht
die historische Interpretation, sondern bildet einen neuen Zustand mit
eigener Authority.

Dasselbe Prinzip gilt für Änderungen innerhalb der Modellbildung.
Model Placement, Model Quality und gegebenenfalls Target-Model Formulation
werden als eigene Entscheidungen persistiert. Ein freigegebener Nachfolger
überschreibt den vorherigen Internal-Model-Zustand nicht rückwirkend.
Dadurch bleibt beispielsweise nachvollziehbar, welche Beschreibung zunächst
in \texttt{IEM-000001} vorlag und welche human-autorisierte Formulierung
anschließend in \texttt{IEM-000002} übernommen wurde. Traceability
beschreibt damit nicht lediglich statische Beziehungen zwischen Artefakten,
sondern auch die Entwicklung autorisierter Zustände über den
Verarbeitungsprozess hinweg.

Auch fehlgeschlagene Verarbeitungsschritte sind Bestandteil dieser
Prozesshistorie. Der in Abschnitt~\ref{sec:results_traceability}
dargestellte fehlgeschlagene Processing Attempt bleibt einschließlich seines
Fehlerzustands und der Raw Response erhalten, obwohl ein späterer Retry
erfolgreich weiterverarbeitet wurde. Der erfolgreiche Zustand ersetzt den
fehlgeschlagenen Attempt nicht. Dadurch kann nicht nur rekonstruiert werden,
welche Information in das Modell gelangt ist, sondern auch, warum ein
bestimmter maschineller Zwischenzustand nicht weiterverwendet wurde.

Diese Form der Traceability unterscheidet sich von einer ausschließlich
nachträglichen Verknüpfung fertiger Engineering-Artefakte. Traceability wird
im entwickelten Architekturkonzept während der Verarbeitung aufgebaut.
Source-, Evidence-, Review-, Authority- und Modellreferenzen werden beim
Entstehen der jeweiligen Zustände persistiert und an nachfolgende
Verarbeitungsschritte weitergegeben. Anforderungen an nachvollziehbare
Beziehungen zwischen Engineering-Information und ihren abgeleiteten
Artefakten sind auch im Requirements Engineering etabliert
\cite{ISO29148}. Der untersuchte Ansatz erweitert diese Logik auf die
probabilistischen und human-autorisierten Zwischenzustände eines
KI-gestützten Transformationsprozesses.

Dabei ist zwischen \emph{Provenance} und \emph{Authority} zu unterscheiden.
Eine Information kann eindeutig auf eine Source zurückgeführt werden, ohne
dadurch bereits fachlich autorisiert zu sein. Umgekehrt besitzt eine Human
Decision nur dann eine nachvollziehbare Grundlage, wenn der Zustand, auf den
sie sich bezieht, eindeutig identifizierbar bleibt. Für den entwickelten
Ansatz werden deshalb beide Dimensionen gemeinsam geführt: Provenance
beschreibt, woher ein Informationszustand stammt, während die
Authority-Kette beschreibt, durch welche Entscheidung er für einen
nachfolgenden Verarbeitungsschritt zugelassen wurde.

Diese Trennung ist ebenfalls für die finale technische Validierung relevant.
Der Validation State ist an einen konkreten generierten Artifact Set
gebunden. Die anschließende Human Release Decision referenziert wiederum
diesen validierten Zustand und dessen Fingerprint. Dadurch wird ein
erfolgreiches Validation-Ergebnis nicht auf beliebige spätere Änderungen des
Modells übertragen. Ein veränderter Modellzustand erfordert einen neuen
Validierungs- und Freigabepfad. Technische Validität und Human Release bleiben
damit auch innerhalb der Traceability unterschiedliche Zustände.

Die formative Evaluation führte an mehreren Stellen zu einer Präzisierung
dieser Traceability-Grenzen. Besonders die Multi-Source-Erprobung zeigte,
dass lokal eindeutige Artefakt-IDs für eine projektweite Verarbeitung nicht
ausreichen, wenn derselbe Identifier in unterschiedlichen Processing Runs
auftreten kann. Die darauf folgende Architekturkorrektur erweiterte die
Identität um den jeweiligen Processing-Kontext und band Project Fit an den
konkreten Source-, Run-, Attempt- und Authority-Zustand. Diese
finding-getriebene Änderung wurde als Architekturentscheidung dokumentiert
und anschließend im Multi-Source-Retest überprüft
(vgl. Anhang~\ref{app:findings_register}).

Aus der Evaluation folgt jedoch keine Aussage über eine vollständige
Traceability sämtlicher in der Realität relevanten Engineering-Information.
Nachvollziehbar ist die Verarbeitung innerhalb der registrierten und für das
Projekt zugelassenen Informationsbasis. Eine Information, die nicht als
Source vorliegt oder nicht in den betrachteten Verarbeitungspfad aufgenommen
wurde, kann durch die Traceability nicht nachträglich rekonstruiert werden.
Die Architektur unterstützt damit die kontrollierte Herleitung innerhalb
einer definierten Informationsbasis, nicht deren fachliche Vollständigkeit.

Ebenso ist der Umfang der persistierten Evidence nicht mit einer
Qualitätsaussage über den Inhalt gleichzusetzen. Eine große Zahl von
Zwischenartefakten, Review Items oder Traceability Entries ist kein
eigenständiger Nachweis für Modellqualität. Dies zeigte bereits die formative
Erprobung, in der die reine Anzahl erzeugter Review Items keine belastbare
Aussage über Source Grounding, Relevance oder tatsächliche
Interpretationsvarianz erlaubte
(vgl. Anhang~\ref{app:findings_register}). Entscheidend ist vielmehr, ob die
für eine konkrete Engineering-Aussage relevanten Vorgänger-, Entscheidungs-
und Authority-Zustände eindeutig miteinander verbunden sind.

Für das entwickelte Architekturkonzept ist Traceability damit nicht als
zusätzliche Dokumentationsfunktion nach Abschluss der Modellgenerierung zu
verstehen. Sie entsteht aus der persistierten Verkettung der kontrollierten
Informations- und Entscheidungszustände selbst. Das resultierende
SysML-v2-Modell bildet den veröffentlichten Engineering-Zustand ab; der
zugehörige Evidence-, Decision-, Provenance- und Traceability-Artefaktsatz
ermöglicht die Rekonstruktion, wie dieser Zustand innerhalb des untersuchten
Verarbeitungspfads zustande gekommen ist.

\subsection{Multi-Source-Verarbeitung und projektweite Integration}
\label{sec:discussion_multisource}

Die Multi-Source-Erprobung zeigt, dass die kontrollierte Verarbeitung mehrerer
Engineering Sources eine zusätzliche Architekturgrenze erfordert. Eine
source-lokale Verarbeitung allein genügt nicht, sobald autorisierte
Informationen aus mehreren Quellen in einen gemeinsamen Modellbildungsprozess
überführt werden sollen. Gleichzeitig darf die projektweite Integration die
Herkunft und den jeweiligen Authority-Zustand dieser Informationen nicht
auflösen. Für das entwickelte Architekturkonzept ergibt sich daraus die
Anforderung, source-lokale Autorisierung und projektweite Synthese explizit
voneinander zu trennen.

Diese Trennung ist für Brownfield-Szenarien besonders relevant. Bestehende
Engineering-Information liegt typischerweise nicht in einem einzelnen
konsistenten Modell vor, sondern verteilt sich auf unterschiedliche
Dokumente, Disziplinen und Repräsentationsformen. Ein wesentliches Ziel von
MBSE besteht darin, solche Informationen stärker über gemeinsame
Modellstrukturen und nachvollziehbare Beziehungen zu integrieren
\cite{Hendriks2022}. Auch SysML-zentrierte Integrationsansätze adressieren
die Zusammenführung unterschiedlicher Engineering-Sichten innerhalb eines
gemeinsamen Systemmodells \cite{Zhao.2023}. Für eine KI-gestützte Migration
entsteht jedoch zusätzlich die Frage, wie die Authority der jeweiligen
Quellinformation während dieser Integration erhalten bleibt.

Der erste formale Multi-Source-Test des Prototyps machte diese Grenze durch
eine Identitätskollision zwischen Processing-Artifacts sichtbar. Lokal
eindeutige Artefaktidentitäten reichten für die projektweite Verarbeitung
nicht aus, weil gleiche lokale IDs in unterschiedlichen Runs auftreten
konnten. Damit war die Zuordnung eines nachgelagerten Artefakts zu seinem
Source- und Processing-Kontext nicht in jedem Fall eindeutig. Der
Verarbeitungspfad wurde an dieser Stelle nicht fortgesetzt. Die historische
Beobachtung, die darauf folgende Architekturkorrektur und der erfolgreiche
Retest sind im Findings-Register dokumentiert
(vgl. Anhang~\ref{app:findings_register}).

Die daraus abgeleitete Architekturkorrektur behandelte dieses Problem nicht
durch eine frühzeitige Zusammenführung der source-lokalen Zustände.
Stattdessen wurde die Identität projektweit um den Processing-Kontext
erweitert. Ein Processing Artifact wird damit nicht ausschließlich durch
seine lokale Artifact-ID bestimmt, sondern durch den Kontext, in dem dieser
Zustand entstanden ist. Für die projektweite Weiterverarbeitung bleibt
dadurch nachvollziehbar, aus welchem Processing Run und damit aus welchem
Source-bezogenen Informationspfad ein Artefakt stammt.

Diese Erkenntnis führte zugleich zur Präzisierung des Project-Fit-Schritts.
Project Fit entscheidet nicht erneut über den fachlichen Inhalt einer bereits
human-autorisierten Source. Der Schritt bindet vielmehr den konkreten
source-lokalen Authority-Zustand an den Projektkontext und entscheidet, ob
dieser Zustand für die gemeinsame Modellbildung zugelassen wird. Im finalen
Verarbeitungspfad entsteht daher die Struktur

\begin{center}
source-lokales Processing\\
$\downarrow$\\
source-lokale Human Authority\\
$\downarrow$\\
Approved Engineering Information\\
$\downarrow$\\
Project Fit / Admission\\
$\downarrow$\\
projektweite Modellbildung.
\end{center}

Eine wesentliche Konsequenz besteht darin, dass die einzelnen
source-lokalen Approved-Engineering-Information-Zustände nicht zu einem neuen
synthetischen Authority-Artefakt verschmolzen werden. Mehrere autorisierte
Sources werden gemeinsam als Eingabe für die Modellbildung verwendet, ihre
jeweilige Provenance und Authority bleiben jedoch getrennt referenzierbar.
Damit kann eine spätere Modelleigenschaft auf mehrere Quellen zurückgehen,
ohne dass die Herkunft der einzelnen Beiträge durch die projektweite
Integration verloren geht.

Gleichzeitig dürfen Source-Grenzen die fachliche Synthese nicht künstlich
begrenzen. Zwei unterschiedliche Engineering Sources können Aussagen über
denselben Systemgegenstand oder denselben modellierungsrelevanten Concern
enthalten. Die Source-Provenance beschreibt in diesem Fall die Herkunft der
Information, nicht zwingend die Grenze der späteren Modellstruktur. Im
finalen Architekturstand werden deshalb mehrere zugelassene source-lokale
Informationszustände in einen gemeinsamen Model Candidate Set überführt. Die
Zielmodellbildung bleibt damit projektweit, während die Provenance der
zugrunde liegenden Engineering Information source-lokal erhalten wird.

Gerade diese Unterscheidung war Gegenstand der finding-getriebenen
Architekturentwicklung. Die Multi-Source-Erprobung führte zunächst zu einer
strengeren source-lokalen Abgrenzung, anschließend jedoch auch zu der
Präzisierung, dass Provenance-Grenzen nicht automatisch zu
Modellierungsgrenzen werden dürfen. Die entsprechenden
Architekturentscheidungen wurden in den nachfolgenden Architecture Decision
Records dokumentiert und in der finalen Architektur konsolidiert
(vgl. Anhang~\ref{app:architecture_decision_records}). Der erfolgreiche
Retest zeigt für den untersuchten Proof-of-Concept-Pfad, dass diese beiden
Anforderungen gemeinsam realisiert werden konnten: source-lokale Authority
blieb erhalten, während die zugelassenen Informationen anschließend in einen
gemeinsamen projektweiten Modellbildungsprozess eingingen.

Die textuelle und formale Modellrepräsentation von SysML~v2 unterstützt diese
projektweite Zusammenführung insofern, als unterschiedliche
Engineering-Inhalte in einem gemeinsamen, referenzierbaren Modellraum
materialisiert werden können \cite{SysMLv2}. Die Entscheidung, welche
Informationen überhaupt in diesen Modellraum eingehen dürfen, wird dadurch
jedoch nicht beantwortet. Diese Verantwortung verbleibt in der vorgelagerten
Informations- und Authority-Architektur.

Der Multi-Source-Test bildet zugleich nur einen kontrollierten Ausschnitt
realer Engineering-Integration ab. Die verwendeten Quellen waren synthetisch,
inhaltlich überschaubar und auf denselben Projektkontext abgestimmt.
Automatische Konfliktauflösung zwischen widersprüchlichen autorisierten
Sources, umfangreiches Change Management, Versionierung sich unabhängig
ändernder Quellen sowie eine allgemeine semantische Reconciliation wurden
nicht vollständig umgesetzt. Entsprechende offene Findings wurden nach
Abschluss des vorgesehenen Proof-of-Concept-Umfangs nicht weiter aufgelöst
und bleiben Gegenstand der Limitationen und des Ausblicks
(vgl. Anhang~\ref{app:findings_register}).

Für das entwickelte Architekturkonzept ergibt sich daraus, dass
Multi-Source-Fähigkeit weder durch eine möglichst frühe Zusammenführung aller
Informationen noch durch dauerhaft getrennte source-spezifische Modelle
erreicht wird. Erforderlich ist vielmehr eine kontrollierte Übergabe zwischen
zwei Ebenen: Source Identity, Provenance und Human Authority bleiben lokal
gebunden; die daraus zugelassenen Engineering-Inhalte können anschließend
projektweit konsolidiert und gemeinsam modelliert werden. Diese
Unterscheidung bildet gemeinsam mit den zuvor diskutierten
Evidence-, Human-Authority- und Modellbildungsgrenzen die Grundlage für die
nachfolgende Beantwortung der Forschungsfragen.

\subsection{Beantwortung der Forschungsfragen}
\label{sec:discussion_research_questions}

Die Forschungsfragen werden im Folgenden auf Grundlage des entwickelten
Architekturkonzepts, seiner prototypischen Realisierung und der in
Kapitel~\ref{sec:results} dargestellten Evaluationsergebnisse beantwortet.
Die Antworten sind dabei im Sinne der gestaltungsorientierten Forschung als
begründete Aussagen über das entwickelte und untersuchte Artefakt zu verstehen.
Sie beanspruchen keine universelle Gültigkeit für beliebige
KI- oder MBSE-Systeme, sondern beschreiben die aus Konzeption, Realisierung
und Evaluation abgeleiteten Gestaltungsprinzipien für den untersuchten
Anwendungsfall \cite{Wieringa2014,ISO42030}.


\subsubsection{Teilforschungsfrage 1: Anforderungen an die Eingangsdaten}
\label{sec:discussion_rq1}

Die erste Teilforschungsfrage lautet:

\begin{quote}
\textit{Welche Anforderungen (Constraints) müssen an die Eingangsdaten gestellt
werden, damit ein KI-Agent zuverlässig arbeiten kann?}
\end{quote}

Die Ergebnisse der Arbeit präzisieren zunächst, was unter einem zuverlässigen
Arbeiten eines KI-Agenten in diesem Kontext verstanden werden kann.
Zuverlässigkeit kann nicht bedeuten, dass ein Large Language Model aus
beliebigen Eingangsdaten jederzeit eine fachlich korrekte Interpretation
erzeugt. Untersuchungen zum Einsatz generativer Modelle zeigen, dass auch
sprachlich plausible Ergebnisse sachlich falsch, unvollständig oder
unzureichend begründet sein können
\cite{Topcu2025,Barke2023,Sergeyuk2025}. Die Architektur kann diese
probabilistische Eigenschaft des Modells nicht beseitigen. Sie kann jedoch
kontrollieren, auf welcher Informationsbasis eine Interpretation erfolgt und
unter welchen Bedingungen deren Ergebnis weiterverwendet werden darf.

Damit verschiebt sich die Antwort gegenüber einer rein formatbezogenen
Betrachtung der Eingangsdaten. Die Evaluation hat nicht gezeigt, dass
Engineering Sources bereits einem einheitlichen Schema, einer bestimmten
Dokumentstruktur oder sogar einer SysML-nahen Form entsprechen müssen.
Heterogene textuelle Bestandsinformationen konnten verarbeitet werden, nachdem
sie in eine kontrollierte Source-Repräsentation überführt worden waren.
Entscheidend waren vielmehr die Bedingungen, unter denen eine Source in den
Engineering Information Flow aufgenommen wurde.

Eine erste notwendige Bedingung ist die eindeutige \emph{Identität} der
Informationsquelle. Eine Source muss innerhalb des Projektkontexts stabil
identifizierbar sein und ihr verarbeiteter Zustand muss auf genau diese Source
zurückgeführt werden können. Bei wiederholter oder paralleler Verarbeitung
muss zusätzlich unterscheidbar bleiben, in welchem Run und Attempt ein
resultierender Informationszustand entstanden ist. Die Multi-Source-Erprobung
zeigte unmittelbar, dass eine nur lokal eindeutige Artefaktidentität hierfür
nicht ausreicht. Ohne den Processing-Kontext wäre die Herkunft späterer
Artefakte projektweit nicht mehr eindeutig rekonstruierbar.

Eine zweite Bedingung betrifft die \emph{Informationsrolle} der Source.
Die formative Evaluation zeigte, dass zusätzliche Instruktionen,
Orchestrierungsinformationen oder Kontextwissen von einem Sprachmodell
inhaltlich genauso verarbeitet werden können wie die eigentliche
Engineering Source. Daraus folgt jedoch nicht, dass diese Informationen
dieselbe fachliche Authority besitzen. Das entwickelte Architekturkonzept
unterscheidet deshalb zwischen Engineering Sources, die fachliche Evidence
für Aussagen über das System liefern können, und Context Sources, die die
Interpretation, Klassifikation oder projektbezogene Einordnung unterstützen.
Zusätzliches Wissen wird damit nicht grundsätzlich ausgeschlossen. Seine
Rolle und damit seine zulässige Wirkung im Engineering-Prozess müssen jedoch
explizit begrenzt sein.

Die dritte Bedingung ist die \emph{Rückführbarkeit fachlicher Aussagen auf
konkrete Source-Inhalte}. Hierfür wurde source-grounded Evidence als eigene
Informationsschicht eingeführt. Eine spätere Interpretation darf nicht nur
einen plausiblen Engineering-Inhalt enthalten, sondern muss auf den
zugrunde liegenden Source-Ausschnitt zurückführbar bleiben. Dies entspricht
der allgemeinen Funktion von Provenance, die Herkunft von Daten und die zu
ihnen führenden Verarbeitungsschritte nachvollziehbar zu machen
\cite{MoreauMissier2013,Simmhan2005}. Im Requirements Engineering bildet
die Nachvollziehbarkeit von Herkunft und Beziehungen zwischen
Engineering-Information ebenfalls ein wesentliches Prinzip
\cite{ISO29148}.

Eine vierte Bedingung betrifft den Umgang mit \emph{fehlender oder nicht
eindeutig ableitbarer Information}. Heterogene Legacy-Daten müssen für eine
kontrollierte KI-Verarbeitung nicht vollständig oder widerspruchsfrei sein.
Dies wäre gerade für Brownfield-Szenarien eine unrealistische Voraussetzung.
Die Architektur muss jedoch zwischen vorhandener, aber mehrdeutiger
Information und nicht vorhandener Information unterscheiden. Eine
mehrdeutige Aussage kann als Unsicherheit oder Open Question weitergeführt
werden. Fehlt dagegen eine hinreichende Source-Evidence, darf aus der
statistischen Plausibilität einer möglichen Ergänzung keine autorisierte
Engineering-Aussage entstehen.

Dieser Punkt grenzt die entwickelte Lösung auch von einer bloßen
Informationsanreicherung ab. Für die Modellmigration ist nicht entscheidend,
dass das LLM möglichst viele vermeintlich vollständige Engineering-Aussagen
produziert. Entscheidend ist, dass erkennbar bleibt, welche Aussagen durch
die zugelassene Informationsbasis gestützt werden. Die im Bereich MBSE
angestrebte stärkere Integration und Vernetzung von Engineering-Information
\cite{Hendriks2022} setzt gerade bei einer Migration aus bestehenden
Informationsbeständen voraus, dass die Herkunft der integrierten Information
nicht verloren geht.

Die fünfte Bedingung besteht deshalb in einer kontrollierten
\emph{Admission} der Informationsbasis. Eine vorhandene Datei ist nicht allein
durch ihre technische Verfügbarkeit Bestandteil des autoritativen
Engineering-Prozesses. Sie wird zunächst registriert, in ihrer Rolle
klassifiziert und einem konkreten Projekt zugeordnet. Bei mehreren Sources
bleiben diese Zustände zunächst source-lokal. Erst nach Verarbeitung und
Human Authority werden die daraus hervorgegangenen Engineering-Zustände für
die projektweite Modellbildung zugelassen.

Damit zeigt die Untersuchung zugleich eine Grenze des Begriffs
``Eingabedatenqualität''. Die Architektur kann gewährleisten, dass die
registrierten und zugelassenen Informationen kontrolliert verarbeitet und
ihre Ableitungen zurückverfolgt werden. Sie kann jedoch nicht feststellen,
ob außerhalb dieser Informationsbasis weitere für das reale System relevante
Dokumente oder Sachverhalte existieren. Die Vollständigkeit der ausgewählten
Informationsbasis bleibt daher eine vorgelagerte Engineering-Verantwortung.

\paragraph{Antwort auf Teilforschungsfrage 1.}

Für eine zuverlässige KI-gestützte Verarbeitung müssen die Eingangsdaten
nicht bereits homogen, vollständig oder modellkonform vorliegen. Sie müssen
jedoch als kontrollierte Informationsbasis behandelbar sein. Dazu gehören
eine eindeutige und stabile Source-Identität, eine explizite Zuordnung ihrer
Informationsrolle, eine nachvollziehbare Provenance, die Rückführbarkeit
abgeleiteter Aussagen auf konkrete Source-Evidence sowie ein expliziter
Umgang mit fehlender und mehrdeutiger Information. Kontextwissen muss von
fachlich evidenzgebenden Sources getrennt bleiben. Informationen, die aus der
zugelassenen Source-Basis nicht ableitbar sind, dürfen nicht stillschweigend
ergänzt werden. Zuverlässigkeit entsteht damit nicht durch die Annahme eines
fehlerfreien KI-Agenten, sondern durch die kontrollierte Begrenzung seiner
Informationsbasis und seiner Wirkung auf nachfolgende Engineering-Zustände.


\subsubsection{Teilforschungsfrage 2: Einbindung des Ingenieurs als Kontrollorgan}
\label{sec:discussion_rq2}

Die zweite Teilforschungsfrage lautet:

\begin{quote}
\textit{Wie wird der Ingenieur als Kontrollorgan in die Architektur eingebunden,
um bei detektierten Mehrdeutigkeiten effektiv einzugreifen?}
\end{quote}

Human-in-the-Loop-Ansätze gehen grundsätzlich davon aus, dass menschliche
Expertise an geeigneten Stellen in automatisierte Analyse- und
Entscheidungsprozesse eingebunden wird
\cite{MosqueiraRey2023,Lazaros2026}. Für den hier untersuchten
Engineering-Prozess reicht es jedoch nicht aus, den Menschen lediglich als
abschließenden Prüfer eines bereits erzeugten SysML-v2-Modells vorzusehen.
Die formative Evaluation zeigte, dass fachlich relevante
Authority-Entscheidungen an mehreren voneinander verschiedenen Stellen des
Informationsflusses auftreten.

Ein Ausgangspunkt der Architektur war die Verwendung mehrerer
Persona-Perspektiven zur Interpretation derselben Engineering-Information.
Multi-Agent-Ansätze zeigen, dass die Gegenüberstellung mehrerer
Modellinstanzen Reasoning und Ergebnisqualität unter bestimmten Bedingungen
verbessern kann \cite{Du2024}. Daraus folgt jedoch nicht, dass die
Mehrheitsmeinung mehrerer Agenten eine Engineering-Entscheidung ersetzen
kann. Auch Arbeiten zur menschlichen Intervention in LLM-basierten
Multi-Agent-Verfahren weisen darauf hin, dass die Rolle menschlicher
Entscheidungskompetenz durch den Einsatz mehrerer Modelle nicht grundsätzlich
entfällt \cite{TriemDing2024}.

Die Untersuchung führte hier zu zwei wichtigen Präzisierungen. Erstens müssen
mehrere Interpretationen auf denselben fachlichen Referenzgegenstand bezogen
sein. Die frühere Persona-first-Verarbeitung erzeugte teilweise voneinander
abweichende Subject-Populationen. Ein Vergleich dieser Populationen vermischte
damit die Frage, \emph{welcher} Engineering-Gegenstand erkannt worden war,
mit der Frage, \emph{wie} dieser Gegenstand interpretiert wird. Durch die
Bildung gemeinsamer Canonical Engineering Subjects vor der
Persona-Interpretation wird die Varianz dagegen auf einen gemeinsamen
Referenzgegenstand bezogen.

Zweitens ist fehlende Persona-Varianz nicht mit fehlender Unsicherheit
gleichzusetzen. Mehrere Perspektiven können dieselbe eindeutige Interpretation
liefern. Sie können jedoch ebenso dieselben zwei Alternativen oder dieselbe
Informationslücke erkennen. In diesem Fall besteht Konsens über eine
Ambiguität. Ein rein numerisches Mehrheits- oder Agreement-Kriterium wäre
deshalb kein hinreichender Mechanismus für die automatische Freigabe.
Consensus und Variance werden im entwickelten Ansatz als Review-Evidence
verwendet, nicht als Ersatz für Engineering Authority.

Die effektive Einbindung des Ingenieurs besteht folglich nicht darin, jede
LLM-Ausgabe vollständig manuell nachzuprüfen. Ein solcher Ansatz würde einen
großen Teil des potenziellen Automatisierungsnutzens wieder aufheben.
Stattdessen werden Human Decisions an den Stellen eingeführt, an denen sich
der Status einer Information qualitativ verändert. Besonders relevant ist
der Übergang von maschineller Interpretation zu autorisierter Engineering
Information. Der Engineer erhält hierfür Source-Evidence,
Persona-Interpretationen, Consensus beziehungsweise Variance und erkannte
Unsicherheiten als Entscheidungsgrundlage. Erst seine Entscheidung erzeugt
den für die weitere Verarbeitung autorisierten Zustand.

Die Evaluation zeigte jedoch, dass diese erste Human Authority allein nicht
ausreicht. Die Autorisierung des \emph{Engineering Meaning} beantwortet noch
nicht zwangsläufig, wie dieser Inhalt im Modell repräsentiert werden soll.
Deshalb wurde eine weitere Authority-Grenze für Model Placement eingeführt.
Kann die Zielrepräsentation nicht deterministisch aus den zugelassenen Regeln
abgeleitet werden, wird die Entscheidung nicht innerhalb der Model Assembly
erneut an ein LLM delegiert. Sie verbleibt als unresolved State oder wird
einer Human Decision zugeführt.

Dasselbe Prinzip gilt für die Target-Model Formulation. Eine bereits
autorisierte Bedeutung und eine akzeptierte Platzierung können noch mehrere
zulässige konkrete Modellformulierungen erlauben. Wo dieser Schritt
Engineering-Entscheidungen enthält, muss auch diese Authority von der
anschließenden deterministischen Serialisierung getrennt bleiben. Am Ende des
Prozesses folgt zusätzlich ein Final Model Review mit Human Release. Dieser
entscheidet über das konkrete validierte Modellartefakt und nicht erneut über
jede einzelne vorgelagerte Source-Aussage.

Die Architektur integriert Human Authority damit \emph{gestuft}. Der Mensch
ist nicht ein einziges Gate am Ende des Prozesses, sondern besitzt klar
definierte Entscheidungsverantwortung an unterschiedlichen semantischen und
modellbezogenen Übergängen. Diese Struktur reduziert zugleich das Risiko,
dass eine spätere technische Stufe implizit fachliche Authority übernimmt,
nur weil sie für die Fortsetzung des Workflows einen eindeutigen Wert
benötigt.

Ein weiterer Bestandteil einer effektiven Human Authority ist die
Persistierung der Entscheidungen. Reviewer, Outcome, zugrunde liegender
Informationszustand und gegebenenfalls Rationale bleiben als eigene
Decision-Artefakte erhalten. Dadurch wird die menschliche Entscheidung nicht
zu einem transienten Klick in einer Benutzeroberfläche, sondern zu einem
nachvollziehbaren Bestandteil der Engineering-Provenance. Dies ist
insbesondere relevant, wenn spätere Modellzustände auf vorherigen
Entscheidungen aufbauen.

Die formative Evaluation zeigt gleichzeitig, dass dieser Ansatz noch keine
optimale Human-Machine Interaction garantiert. Mehrere offene Findings
betreffen den Review-Aufwand, die textlastige Darstellung von
Entscheidungsgrundlagen, die Visualisierung gemeinsamer Ambiguität und die
Anzahl erforderlicher Interaktionen. Der Proof of Concept belegt daher nicht,
dass der realisierte Review-Prozess bereits effizient im Sinne einer
industriellen Benutzerstudie ist. Er zeigt vielmehr, wie die
Entscheidungsautorität architektonisch verortet werden kann. Die Optimierung
der konkreten Interaktion bleibt Gegenstand weiterer Entwicklung.

\paragraph{Antwort auf Teilforschungsfrage 2.}

Der Ingenieur sollte nicht als nachgelagerter Prüfer eines vollständig
KI-generierten Modells eingebunden werden, sondern als gestufte
Engineering- und Model Authority an definierten Übergängen des
Informationsflusses. Mehrere KI-Perspektiven dienen dazu,
Interpretationsalternativen, Konsens und Varianz sichtbar zu machen. Diese
Informationen strukturieren den Review, besitzen jedoch selbst keine
Freigabeautorität. Der Engineer autorisiert zunächst das Engineering Meaning
und greift anschließend dort erneut ein, wo Model Placement oder
Target-Model Formulation nicht eindeutig aus den bereits autorisierten
Zuständen und Regeln ableitbar sind. Nicht aufgelöste Mehrdeutigkeit bleibt
explizit unresolved und darf nicht durch einen technischen Default oder eine
erneute implizite LLM-Entscheidung verdeckt werden. Die menschlichen
Entscheidungen werden persistiert und bleiben damit Bestandteil der
Provenance- und Traceability-Kette. Effektive Human Authority bedeutet
folglich nicht maximale manuelle Kontrolle, sondern gezielte
Entscheidungsautorität an den Stellen, an denen probabilistische
Interpretation in einen autorisierten Engineering-Zustand übergeht.


\subsubsection{Teilforschungsfrage 3: Gestaltung ontologischer Leitplanken}
\label{sec:discussion_rq3}

Die dritte Teilforschungsfrage lautet:

\begin{quote}
\textit{Wie müssen die ontologischen Regeln („Leitplanken“) gestaltet sein,
gegen die die Eingangsdaten geprüft werden, um die Plug \& Play-Fähigkeit der
resultierenden SysML v2 Modelle (Modularität) zu gewährleisten?}
\end{quote}

Die Ergebnisse der Arbeit führen bei dieser Forschungsfrage zu einer
wesentlichen Präzisierung ihrer ursprünglichen Annahme. Ontologische Regeln
sind für die kontrollierte Transformation von Engineering-Information
relevant, können die Modularität und unmittelbare Einsetzbarkeit eines
SysML-v2-Modells jedoch nicht allein gewährleisten.

Ontologien stellen explizite Konzepte, Kategorien und Beziehungen bereit und
können dadurch die semantische Interoperabilität heterogener
Engineering-Information unterstützen. Dies entspricht sowohl allgemeinen
Ansätzen formaler Ontologien als auch ihrer Anwendung im Systems Engineering
und MBSE \cite{ISO21838_2_2021,OrellanaMandrick2019,Lu2022}.
Aktuelle Arbeiten zu domänenspezifischen Engineering-Ontologien und
ontology-based MBSE verfolgen ebenfalls das Ziel, gemeinsame Semantik über
Domänen- und Modellgrenzen hinweg nutzbar zu machen
\cite{Liu2025SysDO,Dong2026,Kulvatunyou2022}.

Für den untersuchten Transformationsprozess zeigte sich jedoch, dass eine
semantisch korrekte Klassifikation noch keine eindeutige
Zielmodellrepräsentation bestimmt. Die Aussage, dass eine Information
beispielsweise eine Anforderung, ein logisches Element, ein Verhalten oder
eine fachliche Beziehung beschreibt, legt ihre Bedeutung ein Stück weit fest.
Sie beantwortet aber nicht zwangsläufig, welches konkrete SysML-v2-Konstrukt
im verwendeten Modellframework gewählt werden muss. Die formative Evaluation
zeigte genau an dieser Grenze mehrere Fälle, in denen fachlich plausible
Engineering Information nicht direkt in eine zulässige Zielrepräsentation
überführt werden konnte.

Die Leitplanken müssen deshalb auf mehreren Ebenen wirken. Auf der
\emph{semantischen Ebene} begrenzen sie die zulässigen Information Types,
Relationship-Semantiken und epistemischen Zustände. Auf der
\emph{architektonischen Ebene} legen sie fest, welche Bereiche des
Zielmodells existieren, welche Verantwortlichkeiten diese besitzen und welche
Arten von Information dort grundsätzlich eingeordnet werden dürfen. Die im
Proof of Concept verwendete RFLP-Struktur stellt hierfür einen
Modellierungsrahmen bereit. Die Trennung von Requirements, Functional,
Logical und Physical Perspectives ist als strukturierendes Prinzip auch in
der MBSE-Literatur etabliert \cite{LiLockettLawson2020}.

Eine dritte Ebene betrifft die \emph{Target-Model Representation}. Hier muss
explizit geregelt sein, welche Engineering Information durch welche
SysML-v2-Konstrukte repräsentiert werden darf. Diese Abbildung kann nicht
vollständig aus einer allgemeinen Ontologie abgeleitet werden, da sie auch
vom gewählten Modellierungsprofil, der Modellierungsabsicht und dem
zugelassenen Scope der Zielarchitektur abhängt. Der Befund, dass
Engineering Meaning und Target Representation getrennte
Verantwortungsbereiche darstellen, war deshalb eine zentrale
Architekturkorrektur der formativen Evaluation.

Eine vierte Ebene bildet die \emph{Target-Model Formulation}. Selbst nach der
Auswahl eines geeigneten Modellkonstrukts muss dessen konkrete
modellsprachliche Ausprägung bestimmt werden. Die formale Sprachdefinition von
SysML~v2 stellt hierfür die zulässigen Sprachmittel bereit
\cite{SysMLv2}. Arbeiten zur automatisierten oder LLM-gestützten
SysML-v2-Erzeugung zeigen, dass eine maschinelle Generierung formal
verarbeitbarer Modellartefakte grundsätzlich realisierbar ist
\cite{Stein2026,Pan.2025,Cibrian.2025}. Die Ergebnisse der vorliegenden
Arbeit zeigen jedoch, dass die technische Generierbarkeit nicht mit der
fachlichen Autorisierung der Repräsentation gleichgesetzt werden darf.

Deshalb folgt als weitere Leitplanke eine
\emph{deterministische Materialisierungsgrenze}. Sobald Engineering Meaning,
Placement und gegebenenfalls Formulation autorisiert sind, darf die
nachfolgende Assembly keine neue semantische Interpretation einführen.
Unvollständige Mappings müssen sichtbar unresolved bleiben, anstatt durch
einen probabilistischen Fallback geschlossen zu werden. Dies war insbesondere
eine Konsequenz aus der formativen Beobachtung, dass eine frühere Model
Assembly für nicht eindeutig abbildbare Relationships erneut
LLM-Perspektiven aufrief und damit ihre vorgesehene Responsibility Boundary
überschritt.

Schließlich muss die resultierende Modellstruktur technisch gegen die
Zielsprache und die verwendete Werkzeugumgebung validiert werden.
Metamodell- und regelbasierte Validierung kann dabei strukturelle und
semantische Konsistenzbedingungen eines SysML-v2-Modells prüfen
\cite{Cibrian.2025b}. Die im Proof of Concept verwendete SYSIDE-Validierung
bildet eine zusätzliche technische Grenze. Ein erfolgreich validiertes
Artefakt ist damit innerhalb des geprüften Sprach- und Werkzeugausschnitts
verarbeitbar. Daraus folgt jedoch noch keine fachliche Korrektheit des
Engineering-Modells.

Auch der Begriff \emph{Plug \& Play} muss deshalb enger gefasst werden.
Eine vollständig kontextunabhängige Einsetzbarkeit beliebiger erzeugter
SysML-v2-Module wurde durch die Arbeit nicht nachgewiesen. Unterstützt werden
kann vielmehr eine kontrollierte Modularität innerhalb eines definierten
Target Frameworks. Voraussetzung dafür sind stabile Identitäten,
klar definierte Modellbereiche, explizite Beziehungen, zulässige
Target Constructs und deterministische Formulierungs- und
Validierungsregeln. Textbasierte Engineering- und
Architecture-as-Code-Ansätze unterstützen eine solche reproduzierbare,
versionierbare und werkzeuggestützte Modellbildung
\cite{Stirbu2022,Bucaioni2025}.

Die Evaluation zeigt gleichzeitig die Grenze dieses Ergebnisses. Die
implementierte Construct-Abdeckung umfasst nur einen definierten
Proof-of-Concept-Scope. Bestimmte SysML-v2-Konstrukte, darunter weitergehende
Ports, Interfaces, Parts, Actions, States und Occurrences, sind noch nicht
allgemein abgedeckt. Eine generelle Plug-\&-Play-Fähigkeit für beliebige
SysML-v2-Modelle kann daher nicht beansprucht werden.

\paragraph{Antwort auf Teilforschungsfrage 3.}

Ontologische Regeln müssen die zulässige Semantik der verarbeiteten
Engineering-Information eindeutig begrenzen, sind für die Modularität der
resultierenden SysML-v2-Strukturen jedoch nur eine Ebene der erforderlichen
Leitplanken. Ergänzend werden Regeln für die Architektur des Zielmodells,
zulässige Placement- und Representation-Entscheidungen, explizite
Target-Model-Formulierungen sowie deterministische Assembly- und
Validierungsgrenzen benötigt. Ontologie beschreibt damit primär,
\emph{was eine Information bedeutet}; die weiteren Leitplanken bestimmen,
\emph{wo und wie sie im Zielmodell repräsentiert werden darf}. Eine
Plug-\&-Play-Fähigkeit kann auf dieser Grundlage innerhalb eines definierten
Target Profiles und der implementierten Construct-Abdeckung unterstützt
werden. Eine allgemeine Modularitätsgarantie allein durch ontologische
Klassifikation lässt sich aus den Ergebnissen dagegen nicht ableiten.


\subsubsection{Hauptforschungsfrage: Gestaltung der Informationsverarbeitungsarchitektur}
\label{sec:discussion_main_rq}

Die Hauptforschungsfrage lautet:

\begin{quote}
\textit{Wie muss eine Informationsverarbeitungsarchitektur gestaltet sein, um
heterogene Bestandsdaten durch ein systematisches Bewertungs- und
Synthese-Framework sicher in valide SysML v2 Strukturen zu überführen?}
\end{quote}

Die drei Teilforschungsfragen zeigen, dass diese Aufgabe nicht durch einen
einzelnen KI-Agenten und auch nicht durch einen direkten
Legacy-to-SysML-v2-Prompt ausreichend beschrieben werden kann. Die zentrale
Gestaltungsaufgabe besteht vielmehr darin, probabilistische
Informationsverarbeitung in einen kontrollierten Engineering-Prozess
einzubetten.

Aktuelle Forschung zeigt, dass Large Language Models grundsätzlich in der
Lage sind, modellierungsnahe beziehungsweise SysML-v2-kompatible Artefakte zu
erzeugen \cite{Stein2026,Pan.2025,Cibrian.2025}. Gleichzeitig zeigen
Untersuchungen zu LLM-basierten Engineering- und Softwareaufgaben, dass
plausible Ausgaben nicht automatisch verlässlich oder fachlich korrekt sind
\cite{Topcu2025,Barke2023,Sergeyuk2025}. Die zentrale Frage ist deshalb nicht
nur, ob ein Modell SysML-v2-Text erzeugen kann, sondern wie nachvollziehbar
begründet wird, welche Engineering-Information in diesen Text gelangen darf.

Die entwickelte Architektur beantwortet diese Frage zunächst durch eine
\emph{kontrollierte Informationsbasis}. Sources werden registriert,
identifiziert und nach ihrer Informationsrolle unterschieden. Engineering
Sources können fachliche Evidence liefern, während Context Sources die
Interpretation unterstützen, ohne selbst unkontrolliert Engineering-Fakten
zu begründen. Aus den zugelassenen Sources wird source-grounded Evidence
gebildet, die konkrete Aussagen an ihre Herkunft bindet. Provenance ist damit
kein nachträgliches Attribut des finalen Modells, sondern Bestandteil der
Informationsverarbeitung von Beginn an. Dies entspricht der grundlegenden
Funktion von Provenance, Entstehung und Verarbeitung von Informationen
rekonstruierbar zu halten \cite{MoreauMissier2013,Simmhan2005}.

Darauf folgt eine \emph{kontrollierte probabilistische Interpretation}.
Canonical Engineering Subjects werden vor der Persona-Interpretation
etabliert. Mehrere Perspektiven beziehen sich dadurch auf denselben
fachlichen Referenzgegenstand. Consensus und Variance dienen zur
Charakterisierung des Interpretationszustands, erhalten aber keine eigene
Engineering Authority. Damit wird die in Multi-Agent-Ansätzen nutzbare
Perspektivenvielfalt \cite{Du2024} mit einer expliziten
Human-in-the-Loop-Struktur verbunden
\cite{MosqueiraRey2023,Lazaros2026,TriemDing2024}.

Der nächste zentrale Bestandteil ist eine \emph{gestufte Human Authority}.
Der Übergang von maschineller Interpretation zu Approved Engineering
Information erfordert eine explizite Engineering Decision. Die daraus
entstehende autorisierte Information bildet die Grundlage der weiteren
Modellbildung. Für Multi-Source-Projekte bleibt diese Authority zunächst
source-lokal. Erst ein expliziter Project-Fit- beziehungsweise
Admission-Schritt erlaubt die Verwendung des konkreten autorisierten
Source-Zustands in der projektweiten Modellbildung. Source-Provenance bleibt
damit erhalten, ohne die spätere Synthese in ein gemeinsames Modell künstlich
an Source-Grenzen zu binden.

Die Modellbildung erfordert anschließend eine weitere
\emph{Trennung von Verantwortlichkeiten}. Die Evaluation zeigte, dass
Engineering Meaning, Model Placement beziehungsweise Representation,
Target-Model Formulation und technische Serialisierung nicht dieselbe
Entscheidung darstellen. Ontologische Regeln und Engineering-Klassifikationen
begrenzen die fachlich zulässigen Bedeutungen
\cite{OrellanaMandrick2019,Lu2022,Dong2026}. Die Zielmodellarchitektur und
SysML-v2-Regeln bestimmen dagegen, wie diese Bedeutung innerhalb eines
konkreten Modellframeworks repräsentiert werden kann
\cite{SysMLv2}. Ist diese Abbildung nicht eindeutig, verbleibt die
Entscheidung bei einer dafür vorgesehenen Human Authority.

Nach diesen Authority-Grenzen wird die Verarbeitung soweit möglich
\emph{deterministisch}. Das Internal Engineering Model materialisiert die
autorisierten Informationen und Modellentscheidungen, ohne deren fachliche
Bedeutung erneut probabilistisch zu verändern. Die anschließende
SysML-v2-Generierung serialisiert einen definierten Modellzustand. Dieses
Prinzip steht in Beziehung zu textbasierten Engineering- und
Architecture-as-Code-Ansätzen, bei denen Artefakte reproduzierbar,
versionierbar und automatisiert erzeugt werden können
\cite{Stirbu2022,Bucaioni2025}.

Ein weiterer Bestandteil der Antwort ist das \emph{fail-closed Verhalten}.
Kann eine Information nicht eindeutig aus der Source abgeleitet, nicht
autorisiert oder nicht mit den zulässigen Modellierungsregeln repräsentiert
werden, darf der Verarbeitungspfad den fehlenden Zustand nicht stillschweigend
ersetzen. Der Zustand bleibt unresolved, die Verarbeitung wird an einer
Contract-Grenze gestoppt oder die Entscheidung wird an eine Human Authority
zurückgegeben. Fail-closed bedeutet dabei nicht, dass der KI-Agent keine
Fehler erzeugen kann. Es bedeutet, dass bestimmte nicht autorisierte oder
strukturell unzulässige Zustände nicht unbemerkt in den veröffentlichten
Modellpfad übergehen dürfen.

Die Architektur lässt sich damit als kontrollierte Folge von
Informations-, Entscheidungs- und Modellzuständen zusammenfassen:

\begin{center}
Engineering Sources\\
$\downarrow$\\
Source-grounded Evidence\\
$\downarrow$\\
Canonical Engineering Subjects\\
$\downarrow$\\
Persona Interpretations und Consensus/Variance\\
$\downarrow$\\
Human Engineering Decisions\\
$\downarrow$\\
Approved Engineering Information\\
$\downarrow$\\
Project Fit und projektweite Admission\\
$\downarrow$\\
Model Candidates und Human Model Placement\\
$\downarrow$\\
Internal Engineering Model\\
$\downarrow$\\
Target-Model Formulation und Model Quality Authority\\
$\downarrow$\\
deterministische SysML-v2-Generierung\\
$\downarrow$\\
technische Validierung\\
$\downarrow$\\
Final Human Model Review und Release\\
$\downarrow$\\
Published Output.
\end{center}

Diese Verarbeitungskette wird durch einen zweiten, nicht weniger wichtigen
Mechanismus ergänzt: die durchgängige Persistierung von Evidence,
Decisions, Provenance und Traceability. Für einen veröffentlichten
Modellzustand bleibt dadurch nicht nur die ursprüngliche Source
identifizierbar. Auch probabilistische Zwischenzustände, menschliche
Entscheidungen, Modellzustände und technische Validierung lassen sich auf
ihre jeweiligen Vorgänger zurückführen. Die in
Abschnitt~\ref{sec:results_traceability} exemplarisch dargestellte
Traceability zeigt, dass für ein finales Modellelement ein Pfad bis zur
zugrunde liegenden Source-Evidence rekonstruiert werden kann. Umgekehrt kann
eine autorisierte Information in Richtung ihrer nachfolgenden
Modellrepräsentation verfolgt werden. Anforderungen an nachvollziehbare
Beziehungen zwischen Engineering-Artefakten sind im Requirements Engineering
grundsätzlich etabliert \cite{ISO29148}; die vorliegende Architektur erweitert
diese Logik auf die probabilistischen und human-autorisierten
Zwischenzustände der KI-gestützten Transformation.

Auch die Begriffe \emph{sicher} und \emph{valide} aus der
Hauptforschungsfrage müssen vor dem Hintergrund der Ergebnisse präzisiert
werden. ``Sicher'' bezeichnet im Kontext dieser Arbeit keinen Nachweis
funktionaler Sicherheit und keine Safety-Zertifizierung. Gemeint ist eine
kontrollierte Verarbeitung, bei der Source-, Provenance- und
Authority-Grenzen explizit sind und nicht zulässige Zustände fail-closed
behandelt werden. Ebenso bezeichnet ein ``valides SysML-v2-Modell'' zunächst
ein Artefakt, das die definierten internen Regeln erfüllt und von der
verwendeten SysML-v2-Werkzeugkette technisch verarbeitet werden kann.
SysML~v2 stellt dafür eine formal definierte Zielnotation bereit
\cite{SysMLv2}; technische Validierung und semantische Konsistenzprüfung sind
jedoch von der fachlichen Engineering-Korrektheit zu unterscheiden
\cite{Cibrian.2025b}.

Diese Differenzierung erklärt auch die Notwendigkeit des abschließenden Human
Model Review. Ein technisch valides Modell kann syntaktisch und strukturell
korrekt sein und dennoch eine fachlich ungeeignete Aussage enthalten.
Technische Validation beantwortet deshalb die Frage, ob das konkrete
Artefakt die geprüften technischen Bedingungen erfüllt. Die finale Human
Release entscheidet dagegen, ob genau dieser validierte Modellzustand als
Engineering-Output veröffentlicht werden darf.

Die Ergebnisse des Referenzlaufs zeigen, dass die vollständige Kette im
implementierten Proof-of-Concept-Scope ausgeführt werden konnte. Gemeinsam
mit der formativen Evaluation und dem Multi-Source-Retest unterstützt dies
die Schlussfolgerung, dass die einzelnen Architekturgrenzen nicht nur
konzeptionell beschrieben, sondern in einem verbundenen Verarbeitungspfad
prototypisch realisierbar sind. Daraus folgt jedoch weder eine allgemeine
Garantie fachlich korrekter Modellgenerierung noch eine Aussage über
Produktionsreife. Die Evaluation bezieht sich auf die registrierten und
zugelassenen Projektinformationen, den implementierten SysML-v2-Ausschnitt
und die untersuchten Testfälle.

\paragraph{Antwort auf die Hauptforschungsfrage.}

Eine Informationsverarbeitungsarchitektur zur KI-gestützten Überführung
heterogener Bestandsdaten in SysML~v2 muss die probabilistische
KI-Verarbeitung von den autoritativen Engineering- und
Modellentscheidungen trennen. Erforderlich sind eine kontrollierte
Informationsbasis mit expliziten Source-Rollen, source-grounded Evidence und
durchgängiger Provenance, eine gemeinsame Subject Identity vor
Multi-Perspektiven-Interpretation, gestufte Human Authority, eine explizite
Trennung von Engineering Meaning, Model Placement, Target-Model Formulation
und technischer Serialisierung sowie deterministische und fail-closed
kontrollierte Übergänge nach den jeweiligen Authority-Grenzen. Bei
Multi-Source-Verarbeitung müssen Source Identity und Authority lokal erhalten
bleiben, während autorisierte Informationen nach expliziter Admission
projektweit synthetisiert werden können. Der erzeugte SysML-v2-Zustand wird
anschließend technisch validiert und erst durch eine getrennte Human Release
zur Veröffentlichung autorisiert. Parallel dazu muss ein persistierter
Evidence-, Decision-, Provenance- und Traceability-Artefaktsatz erhalten
bleiben, durch den die Herleitung des veröffentlichten Modellzustands
rekonstruiert werden kann.

Die zentrale Antwort der Arbeit besteht damit nicht in einem Verfahren, das
die Unsicherheit probabilistischer KI eliminiert. Sie besteht in einer
Architektur, die diese Unsicherheit von Engineering Authority trennt,
sichtbar macht und ihre Wirkung auf das resultierende Modell kontrolliert.
Das wissenschaftliche Ergebnis ist folglich nicht primär die automatisierte
Erzeugung von SysML-v2-Code, sondern die Gestaltung eines kontrollierbaren
Engineering-Prozesses, innerhalb dessen KI-generierte Interpretationen in
nachvollziehbare, human-autorisierte und technisch validierte Modellzustände
überführt werden können.

\subsection{Limitationen}
\label{sec:discussion_limitations}

Die Ergebnisse dieser Arbeit sind vor dem Hintergrund des gewählten
Forschungsdesigns und des Umfangs der prototypischen Realisierung zu
interpretieren. Ziel der Untersuchung war die Entwicklung und Evaluation
eines Architekturkonzepts für eine kontrollierte KI-gestützte Transformation
heterogener Engineering-Information in SysML~v2. Gegenstand der Evaluation
war damit primär die Frage, ob die entwickelten Informations-,
Verantwortungs- und Authority-Grenzen in einem verbundenen
Proof-of-Concept-Pfad realisierbar und überprüfbar sind. Eine umfassende
Bewertung eines produktionsreifen Engineering-Werkzeugs war nicht Gegenstand
der Arbeit. Diese Abgrenzung entspricht dem gestaltungsorientierten
Charakter der Untersuchung, bei dem ein Artefakt in Bezug auf definierte
Ziele und Evaluationsfragen untersucht wird
\cite{Wieringa2014,ISO42030}.

\paragraph{Datenbasis und empirische Breite}

Die für die formative Evaluation und den abschließenden Referenzlauf
verwendeten Testdaten basieren auf realen Engineering-Inhalten von PreciPoint.
Sie wurden jedoch nicht als unveränderte Produktivdokumente in die
Untersuchung übernommen. Umfang, sprachliche Form und Dokumentstruktur wurden
gezielt reduziert und vereinheitlicht, damit relevante fachliche Aussagen
eindeutig abgegrenzt und die daraus entstehenden Informationspfade
reproduzierbar untersucht werden konnten. Die verwendeten Dokumente und ihre
fachlichen Inhalte sind in
Anhang~\ref{app:vv_test_data} beziehungsweise für den Referenzlauf zusätzlich
in Anhang~\ref{app:traceability_930444_inputs} dokumentiert.

Diese kontrollierte Aufbereitung erhöht die Nachvollziehbarkeit der einzelnen
Testfälle, begrenzt jedoch zugleich die empirische Breite der Evaluation.
Reale industrielle Datenbestände können wesentlich umfangreicher,
uneinheitlicher und stärker miteinander verflochten sein. Unterschiedliche
Dokumentversionen, umfangreiche Tabellen, bestehende Modelle,
disziplinübergreifende Artefakte oder widersprüchliche historische
Informationsstände wurden nicht in einer groß angelegten industriellen
Fallstudie untersucht. Der erfolgreiche Referenzlauf und der
Multi-Source-Retest belegen daher die Ausführbarkeit des entwickelten
Architekturkonzepts innerhalb des untersuchten Scopes, erlauben jedoch keine
statistische Aussage über dessen Robustheit bei beliebigen industriellen
Datenbeständen.

Die bewusste Reduktion der Testdokumente hatte zusätzlich einen praktischen
Grund. Wiederholte LLM-Verarbeitung über mehrere Persona- und
Verarbeitungsstufen verursacht mit wachsendem Informationsumfang einen
zunehmenden Verarbeitungsaufwand. Die Testdaten wurden deshalb so gewählt,
dass sie hinreichend unterschiedliche Engineering-Information für die
Untersuchung der Architekturgrenzen enthielten, ohne die einzelnen
Evaluationsläufe unnötig durch redundante Inhalte zu vergrößern.

\paragraph{Keine unabhängige Evaluation}

Die formative Evaluation und die Human Reviews wurden durch den Entwickler
und Autor der Arbeit durchgeführt. Dies ermöglichte eine enge Rückkopplung
zwischen beobachteten Findings, Architekturkorrekturen und Retests, stellt
jedoch zugleich eine methodische Begrenzung dar. Die Bewertung ist damit
nicht unabhängig von der Entwicklung des untersuchten Artefakts.

Insbesondere die beobachtete Nutzbarkeit der Review-Mechanismen kann deshalb
nicht als Ergebnis einer Usability- oder Anwenderstudie interpretiert werden.
Die Evaluation zeigt, dass die vorgesehenen Human-Authority-Schritte
technisch und prozessual ausgeführt werden konnten. Sie zeigt nicht, wie
effizient, verständlich oder akzeptabel diese Interaktionen für eine größere
Gruppe von Systems Engineers im praktischen Einsatz wären. Die im
Findings-Register dokumentierten Beobachtungen zu textlastigen Reviews,
Interaktionsaufwand und Statusdarstellung unterstreichen diese Grenze
(vgl. Anhang~\ref{app:findings_register}).

\paragraph{Begrenzte Aussage zur Qualität der KI-Interpretation}

Die Arbeit untersucht keine allgemeine Leistungsfähigkeit eines bestimmten
Large Language Models. Weder Accuracy, Precision und Recall der
Engineering-Information-Extraktion noch die Qualität verschiedener
LLM-Anbieter oder Promptstrategien wurden systematisch vergleichend
quantifiziert. Die Anzahl erfolgreicher Verarbeitungsschritte,
Persona-Übereinstimmungen oder verarbeiteter Tokens stellt daher keinen
Qualitätsindikator des verwendeten Sprachmodells dar.

Generative Modelle können plausible, aber sachlich unzutreffende oder
unzureichend begründete Ergebnisse erzeugen
\cite{Topcu2025,Barke2023,Sergeyuk2025}. Das entwickelte Architekturkonzept
versucht deshalb nicht, eine fehlerfreie LLM-Interpretation nachzuweisen.
Untersucht wurde vielmehr, wie deren Ergebnisse source-grounded,
überprüfbar und durch explizite Authority-Grenzen kontrolliert
weiterverarbeitet werden können. Aussagen über die absolute semantische
Leistungsfähigkeit des eingesetzten Modells liegen außerhalb des
Untersuchungsumfangs.

\paragraph{Informationsbasis und fachliche Vollständigkeit}

Die Architektur kann die Verarbeitung registrierter und zugelassener Sources
kontrollieren, nicht jedoch die Vollständigkeit der zugrunde liegenden
Engineering-Informationsbasis garantieren. Ist eine relevante Information
nicht in den aufgenommenen Quellen enthalten oder wurde eine relevante
Source nicht registriert, kann der Verarbeitungspfad diese Information nicht
rekonstruieren.

Source-grounded Evidence und Provenance ermöglichen damit die
Unterscheidung zwischen tatsächlich vorhandener Evidence, daraus
abgeleiteter Interpretation und nicht belegter Information. Sie gewährleisten
jedoch nicht, dass sämtliche für das reale System relevanten Informationen
Teil des Ausgangsdatenbestands sind. Die Auswahl geeigneter Engineering
Sources und die Beurteilung ihrer fachlichen Vollständigkeit bleiben
vorgelagerte Engineering-Aufgaben.

\paragraph{Begrenzter SysML-v2- und Target-Model-Scope}

Die prototypische Realisierung unterstützt nur einen definierten Ausschnitt
der in SysML~v2 verfügbaren Modellierungsmöglichkeiten. Die
Sprachspezifikation bietet einen wesentlich größeren Umfang an Element-,
Relationship-, Behavior- und Strukturkonzepten \cite{SysMLv2}. Im Proof of
Concept wurden gezielt diejenigen Konstrukte implementiert, die für die
Untersuchung des kontrollierten Informationsflusses erforderlich waren.

Das Findings-Register weist weiterhin eine unvollständige Abdeckung unter
anderem weitergehender Ports, Interfaces, Parts, Actions, States und
Occurrences aus
(vgl. Anhang~\ref{app:findings_register}). Auch die
Target-Model-Formulation ist deshalb nicht für beliebige
SysML-v2-Konstrukte generalisiert. Die erfolgreiche technische Validierung
des Referenzmodells belegt folglich die Verarbeitbarkeit des implementierten
Modellierungsausschnitts und keine vollständige Abdeckung der Sprache.

Eng damit verbunden ist die Verwendung eines festgelegten
Modellierungsframeworks. Die im Prototyp verwendete Struktur begrenzt die
zulässigen Placement- und Representation-Entscheidungen und erleichtert
dadurch die kontrollierte Modellbildung. Ob dieselben Leitplanken ohne
Anpassung auf andere Architekturframeworks, Modellierungsprofile oder
domänenspezifische SysML-v2-Methoden übertragen werden können, wurde nicht
untersucht.

\paragraph{Grenzen der semantischen Konsolidierung und Engineering Evolution}

Die Untersuchung zeigt, dass mehrere source-lokale Informationszustände
unter Erhalt ihrer Provenance in einen gemeinsamen projektweiten
Modellbildungsprozess überführt werden können. Eine allgemeine automatische
semantische Reconciliation wurde jedoch nicht realisiert. Insbesondere
können semantisch ähnliche Relationships mit unterschiedlichen Prädikaten
weiterhin getrennt bestehen. Änderungen bereits autorisierter Subjects
können zudem Folgeentscheidungen über Relationships oder nachgelagerte
Modellzustände erforderlich machen.

Die Architektur besitzt bereits explizite Identitäten, Revisionen und
Vorgängerbeziehungen, bildet jedoch noch keinen vollständigen automatischen
Change-Impact-Prozess ab. Abhängige Entscheidungen werden nach einer
fachlichen Änderung nicht generell selektiv invalidiert und automatisch neu
berechnet. Im gegenwärtigen Proof of Concept kann hierfür ein erneuter
fokussierter Verarbeitungslauf erforderlich sein. Diese Grenzen sind im
Findings-Register dokumentiert und wurden nach Abschluss des vorgesehenen
Evaluationsumfangs nicht weiter aufgelöst.

\paragraph{Effizienz und Tokenbedarf}

Der Referenzlauf~930444 benötigte für die LLM-gestützten
Verarbeitungsschritte insgesamt 35 eindeutige LLM-Aufrufe mit 69\,507 Input
Tokens und 21\,693 Output Tokens. Der Gesamtumfang von 91\,200 Tokens umfasst
auch die während des Laufs aufgetretenen fehlgeschlagenen Processing Attempts
und ihre nachfolgenden Retries
(vgl. Tabelle~\ref{tab:results_930444_tokens}).

Die Werte zeigen, dass die mehrstufige Verarbeitung mit mehreren
Interpretationsperspektiven einen nicht vernachlässigbaren
Verarbeitungsaufwand erzeugt. Bereits die formative Evaluation zeigte bei
größer gefassten, source-weiten Verarbeitungsschritten eine starke Zunahme
der erforderlichen LLM-Aufrufe. Dies war ein Grund, spätere
Verarbeitungsschritte stärker auf abgegrenzte Evidence- und Subject-Einheiten
zu beziehen. Eine systematische Optimierung von Tokenverbrauch, Laufzeit oder
Kosten war jedoch nicht Gegenstand der Arbeit. Die gemessenen Werte sind
daher als deskriptiver Referenzwert für den untersuchten Lauf und nicht als
allgemeine Effizienzaussage zu verstehen.

\paragraph{Human Authority und Skalierbarkeit}

Die explizite Human Authority ist ein zentrales Gestaltungsprinzip der
entwickelten Architektur, kann jedoch selbst einen Skalierungsengpass
darstellen. Mit wachsender Anzahl von Sources, Engineering Subjects,
Relationships und Modellentscheidungen wächst auch die Zahl möglicher
Review-Entscheidungen. Human-in-the-Loop-Ansätze müssen grundsätzlich einen
geeigneten Ausgleich zwischen Automatisierung und menschlicher Kontrolle
finden \cite{MosqueiraRey2023,Lazaros2026}.

Die Arbeit zeigt, wie Consensus, Variance und unresolved States für eine
gezieltere Review-Zuführung genutzt werden können. Sie zeigt jedoch nicht,
dass der daraus resultierende Review-Aufwand für große industrielle Projekte
bereits ausreichend gering ist. Die formative Evaluation weist vielmehr auf
Optimierungspotenzial bei Review-Priorisierung, Darstellung und
Interaktionsaufwand hin. Eine quantitative Untersuchung des tatsächlichen
Zeitgewinns gegenüber manueller Modellierung wurde nicht durchgeführt.

\paragraph{Technische Validierung und fachliche Korrektheit}

Die technische SysML-v2-Validierung stellt eine weitere bewusste
Ergebnisgrenze dar. Ein erfolgreicher SYSIDE-Lauf zeigt, dass das erzeugte
Artefakt innerhalb der untersuchten Sprach- und Werkzeugumgebung technisch
verarbeitet werden kann. Daraus folgt weder, dass jede im Modell enthaltene
Engineering-Aussage fachlich korrekt ist, noch dass das Modell für einen
bestimmten späteren Analyse- oder Entwicklungszweck vollständig geeignet ist.

Aus diesem Grund bleibt die technische Validation im Architekturkonzept von
Engineering Authority und finaler Human Model Authority getrennt. Die
Evaluation bestätigt diese Trennung im implementierten Prozess, ersetzt aber
keine domänenspezifische fachliche Validierung durch unabhängige
Fachexperten.

\paragraph{Produktionsreife und technische Systemgrenzen}

Der Turing Generator wurde als Proof of Concept zur Untersuchung des
Architekturkonzepts entwickelt. Aspekte eines produktiven Systems wie
umfassende Benutzer- und Rollenverwaltung, Security, skalierbare
Hintergrundverarbeitung, robuste Job-Steuerung, Monitoring,
Betriebsverfügbarkeit und umfassende Fehlerwiederherstellung standen nicht im
Zentrum der Arbeit. Mehrere im Findings-Register verbliebene Integrations-
und UX-Beobachtungen betreffen genau diese Ebene.

Diese Punkte wurden nach erfolgreichem Abschluss des definierten
Proof-of-Concept-Pfads bewusst nicht weiter optimiert, sofern sie die
Untersuchung der Architektur nicht mehr blockierten. Architekturwirksame
Findings wurden dagegen in Architecture Decisions überführt, implementiert
und erneut geprüft. Die verbleibenden nicht blockierenden Findings markieren
damit die Grenze zwischen dem untersuchten Forschungsartefakt und einer
möglichen Produktweiterentwicklung.

Zusammenfassend begrenzen diese Aspekte insbesondere die
Generalisierbarkeit der Ergebnisse. Gezeigt werden konnte ein kontrollierter
End-to-End-Pfad auf einer aus realen Engineering-Inhalten abgeleiteten,
gezielt aufbereiteten Informationsbasis mit source-grounded
Informationsverarbeitung, persistierter Provenance, gestufter Human
Authority, projektweiter Multi-Source-Integration, deterministischer
Modellbildung und technischer SysML-v2-Validierung innerhalb des definierten
Scopes. Nicht gezeigt wurden eine universelle semantische Korrektheit,
vollständige SysML-v2-Abdeckung, industrielle Skalierbarkeit oder
Produktionsreife. Die daraus entstehenden Weiterentwicklungsbedarfe werden
im Ausblick aufgegriffen.
